% Irodov user-supplied Arihant PDF, printed pp.35--36, PDF pp.44--45.
% Final subsection 1.4: terrestrial and lunar satellites and cosmic velocities.
\exproblem{221}{35}
在地球极点使一物体获得竖直向上的初速 $v_0$，
忽略空气阻力。若已知地球半径 $R$ 与地表重力加速度
$g$，求物体可上升的最大高度 $h$。
\exsolution
取地球为球对称天体，记引力参数
$\mu=GM_E=gR^2$。
地表初态的单位质量机械能为
$\varepsilon_i=v_0^2/2-\mu/R$。
如果初速小于逃逸速度，最高点速率为零、半径为 $R+h$，
故机械能守恒：
\[
\frac{v_0^2}{2}-\frac{\mu}{R}
=-\frac{\mu}{R+h}.
\]
移项可得
\[
\boxed{h=\frac{Rv_0^2}{2gR-v_0^2}}.
\]
\excheck 当 $v_0\to0$，$h\simeq v_0^2/(2g)$；
当 $v_0\to\sqrt{2gR}^{\,-}$，$h\to\infty$。
若 $v_0\ge\sqrt{2gR}$，物体在此理想模型中不再有有限最高点。

\exproblem{222}{35}
一颗人造卫星绕地球作圆周运动。
相对不随地球自转的惯性参考系，它的圆轨道速率为 $v$。
给定地球半径 $R$ 和地表重力加速度 $g$，
求卫星离地高度 $h$。
\exsolution
卫星轨道半径为 $r=R+h$，
由向心运动所需加速度等于引力加速度：
\[
\frac{v^2}{r}=\frac{GM_E}{r^2}
=\frac{gR^2}{r^2}.
\]
所以 $r=gR^2/v^2$，从而
\[
\boxed{h=R\left(\frac{gR}{v^2}-1\right)}.
\]
\excheck 近地表圆轨道对应 $v=\sqrt{gR}$、$h=0$；
高空卫星速率较小。这里的 $v$ 不是相对自转地面的速率。

\exproblem{223}{35}
求在地球赤道上空相对地面静止的
\emph{地球同步静止卫星}轨道半径，
以及它相对于地心惯性系的速率和向心加速度。
\exsolution
卫星必须在赤道平面内与地球同向转动，
其角速度等于地球自转角速度
$\Omega_E=2\pi/T_{\rm sidereal}$，
其中恒星日 $T_{\rm sidereal}\simeq86164\,\mathrm s$。
圆轨道条件
\[
\Omega_E^2r=\frac{GM_E}{r^2},
\]
因此
\[
\boxed{r=\left(\frac{GM_E}{\Omega_E^2}\right)^{1/3}
\simeq4.2164\times10^4\,\mathrm{km}},
\]
\[
\boxed{v=\Omega_Er\simeq3.07\,\mathrm{km/s}},\qquad
\boxed{a=\Omega_E^2r\simeq0.224\,\mathrm{m/s^2}}.
\]
\excheck 给出的 $r$ 是从\emph{地心}量起的半径，
不是离地高度（离地高度约 $3.58\times10^4$ km）；
使用的是恒星日而非恰好24小时的太阳日。

\exproblem{224}{35--36}
一颗卫星在赤道平面内沿半径
$r=2.00\times10^4\,\mathrm{km}$ 的圆轨道自西向东运动。
地面赤道某点每隔 $\tau=11.6\,\mathrm h$
看到卫星再次从头顶经过。
已知引力常数 $G$，求地球质量 $M_E$。
\exsolution
卫星自西向东的角速度 $\omega_s$ 快于地球角速度
$\Omega_E=2\pi/T_E$，因此重复过顶的相对角速度
$\omega_s-\Omega_E=2\pi/\tau$。
故
\[
\omega_s=2\pi\left(\frac1{T_E}+\frac1\tau\right).
\]
根据圆轨道引力与向心条件
$GM_E/r^2=r\omega_s^2$，
\[
\boxed{M_E=\frac{4\pi^2r^3}{G}
\left(\frac1{T_E}+\frac1\tau\right)^2}.
\]
代入 $r=2.00\times10^7\,\mathrm m$、
$\tau=11.6\times3600\,\mathrm s$、
$T_E\simeq86164\,\mathrm s$，
得到 $\boxed{M_E\simeq6.0\times10^{24}\,\mathrm{kg}}$。
\excheck 过顶间隔不是卫星本身的公转周期，
必须减去地球自转角速度后再处理。

\exproblem{225}{36}
一颗卫星沿赤道圆轨道从东向西绕地球运动，
轨道半径 $r=1.00\times10^4\,\mathrm{km}$。
求其在\emph{随地球自转的参考系}中的速率与
加速度大小。
\exsolution
记地心惯性系中的西向圆轨道角速度大小为
$\omega_s=\sqrt{GM_E/r^3}$，
地球自转角速度为 $\Omega_E$（东向）。
在地球旋转参考系中，两者角速度方向相反，
相对西向角速度大小为
$\omega_{\rm rel}=\omega_s+\Omega_E$。
故
\[
\boxed{v_{\rm rel}=r(\omega_s+\Omega_E)
=\sqrt{\frac{GM_E}{r}}+\Omega_Er},
\]
\[
\boxed{a_{\rm rel}=r(\omega_s+\Omega_E)^2}.
\]
取 $GM_E=3.986\times10^{14}\,\mathrm{m^3/s^2}$、
$\Omega_E=7.292\times10^{-5}\,\mathrm s^{-1}$，
得 $v_{\rm rel}\simeq7.04\,\mathrm{km/s}$、
$a_{\rm rel}\simeq4.96\,\mathrm{m/s^2}$。
\excheck 这是旋转坐标系中的\emph{运动学相对加速度}，
并非只有地球引力造成的地心惯性系向心加速度。

\exproblem{226}{36}
两颗卫星在地球赤道附近、紧贴地表的理想圆轨道上运动：
一颗顺地球自转方向，另一颗逆向。
求在\emph{地面旋转参考系}中，
逆向卫星动能是顺向卫星动能的多少倍。
\exsolution
记近地圆轨道地心惯性系速率
$v_c=\sqrt{gR}$、地面赤道线速度
$u_E=\Omega_E R$。
相对地面，顺向与逆向卫星的速率分别为
$v_+=v_c-u_E$、$v_-=v_c+u_E$。
相同卫星质量的动能比即速率平方比：
\[
\boxed{\frac{K_-}{K_+}
=\left(\frac{\sqrt{gR}+\Omega_E R}
{\sqrt{gR}-\Omega_E R}\right)^2
\simeq1.27}.
\]
\excheck 若在地心惯性系，两颗相同半径圆轨道卫星的
动能相等；题目特意要求在地面旋转系比较，
因此出现不同的动能数值。

\exproblem{227}{36}
一颗质量为 $m$ 的月球人造卫星沿接近圆形轨道运动，
初始轨道半径是月球半径 $R_M$ 的 $\eta$ 倍。
卫星受宇宙尘埃的微弱阻力，大小为
$F_d=\alpha v^2$，与速度反向。
若下降非常缓慢、轨道始终可近似视为圆，
求卫星从初始轨道降到月球表面的时间。
\exsolution
记月球引力参数 $\mu=GM_M=g_MR_M^2$。
瞬时圆轨道速率 $v=\sqrt{\mu/r}$，
总机械能为 $E=-\mu m/(2r)$。
阻力耗散功率
\[
-\frac{\mathrm dE}{\mathrm dt}
=F_dv=\alpha v^3=\alpha(\mu/r)^{3/2}.
\]
由 $\mathrm dE/\mathrm dr=\mu m/(2r^2)$，
得到
\[
\frac{\mu m}{2r^2}\frac{\mathrm dr}{\mathrm dt}
=-\alpha\frac{\mu^{3/2}}{r^{3/2}},
\quad
\frac{\mathrm dr}{\mathrm dt}
=-\frac{2\alpha\sqrt\mu}{m}\sqrt r.
\]
从 $r(0)=\eta R_M$ 积分到 $R_M$：
\[
\boxed{t\simeq
\frac{m(\sqrt\eta-1)}{\alpha}
\sqrt{\frac{R_M}{\mu}}
=\frac{m(\sqrt\eta-1)}{\alpha\sqrt{g_MR_M}}}.
\]
\excheck 这是绝热缓慢圆轨道衰减的近似式；
若阻力大到单圈内轨道已明显非圆，需解完整运动方程。

\exproblem{228}{36}
计算月球表面附近的第一宇宙速度（圆轨道速率）
和第二宇宙速度（逃逸速率），并与地球对应数值比较。
\exsolution
对于质量 $M$、半径 $R$ 的球对称天体，
第一宇宙速度来自圆轨道条件
$v_1^2/R=GM/R^2$，故
$v_1=\sqrt{GM/R}=\sqrt{gR}$。
第二宇宙速度来自逃逸总机械能为零：
$v_2=\sqrt{2GM/R}=\sqrt2\,v_1$。
月球取 $g_M\simeq1.62\,\mathrm{m/s^2}$、
$R_M\simeq1.737\times10^6\,\mathrm m$，得
\[
\boxed{v_{1,M}\simeq1.68\,\mathrm{km/s},\qquad
v_{2,M}\simeq2.37\,\mathrm{km/s}}.
\]
相应地球理想值约为
$7.9\,\mathrm{km/s}$ 和 $11.2\,\mathrm{km/s}$。
\excheck 两组速度均满足 $v_2/v_1=\sqrt2$；
这里忽略了大气与天体自转。

\exproblem{229}{36}
一飞船沿近似掠过月球表面的抛物线轨道接近月球。
在最近点瞬时开启制动发动机，使轨道转为掠月球表面的
圆轨道。求制动造成的飞船速率\emph{有符号变化}。
\exsolution
抛物线轨道在无穷远处比机械能为零，
最近点 $r\simeq R_M$ 的速率为
$v_p=\sqrt{2GM_M/R_M}$。
圆轨道所需速率为
$v_c=\sqrt{GM_M/R_M}$。
因此发动机制动造成的速率变化
\[
\boxed{\Delta v=v_c-v_p
=(1-\sqrt2)\sqrt{\frac{GM_M}{R_M}}
\simeq-0.69\,\mathrm{km/s}}.
\]
\excheck 负号表示减速；若题目问制动速度增量的大小，
应报告 $|\Delta v|\simeq0.69\,\mathrm{km/s}$。
不是飞船速率“增加”。

\exproblem{230}{36}
飞船已在地球表面附近的理想圆轨道上运行。
要摆脱地球引力束缚，沿原速度方向最少还需给予多大的
瞬时速率增量？
\exsolution
近地圆轨道速率为
$v_c=\sqrt{GM_E/R_E}=\sqrt{gR_E}$。
在相同高度达到逃逸抛物线轨道需要速率
$v_e=\sqrt{2GM_E/R_E}=\sqrt2\,v_c$。
因此同向瞬时加速的最小速度增量为
\[
\boxed{\Delta v=(\sqrt2-1)\sqrt{gR_E}
\simeq3.27\,\mathrm{km/s}}.
\]
\excheck 这里未包括发动机损耗、轨道变轨时间、
空气阻力及其他天体引力；理想情况下同向脉冲最节省增量。

\exproblem{231}{36}
已知地球质量是月球质量的 $\eta=81$ 倍，
地月中心距离是地球半径 $R_E$ 的 $n=60$ 倍。
求地月连线上两者引力场强合成为零的点距月球中心多远。
\exsolution
记地月中心距 $D=nR_E$，
零场点距月球为 $x$，则距地球为 $D-x$。
两引力方向相反，大小相等：
\[
\frac{GM_M}{x^2}=\frac{GM_E}{(D-x)^2}
=\frac{G\eta M_M}{(D-x)^2}.
\]
取正平方根：
$(D-x)/x=\sqrt\eta$，解得
\[
\boxed{x=\frac{D}{1+\sqrt\eta}
=\frac{nR_E}{1+\sqrt\eta}
=6R_E\simeq3.82\times10^4\,\mathrm{km}}.
\]
\excheck 这是静止牛顿引力场\emph{零场点}，
不是地月旋转参考系中的完整拉格朗日 $L_1$ 点；
后者还需考虑离心惯性项。

\exproblem{232}{36}
一质量为 $m=2.0\times10^3\,\mathrm{kg}$ 的飞船
从地球表面被运至月球表面。
忽略大气、太阳作用和能量损耗，估计所需的最小外力功。
\exsolution
为估算下限，令飞船始末两端相对各自行星静止，
采用地月近似不动、距离 $D$ 固定的保守引力模型。
初态与末态引力势能近似为
\[
U_i\simeq-\frac{GM_E m}{R_E}
-\frac{GM_Mm}{D-R_E},\quad
U_f\simeq-\frac{GM_E m}{D-R_M}
-\frac{GM_Mm}{R_M}.
\]
理想可逆慢速运输时外力净功等于势能差
$W_{\rm ext,min}=U_f-U_i$。
因为 $D\gg R_E,R_M$，忽略远端引力修正得到
\[
W_{\rm ext,min}\simeq
m\left(\frac{GM_E}{R_E}-\frac{GM_M}{R_M}\right)
=m(g_ER_E-g_MR_M).
\]
用 $g_E\simeq9.81\,\mathrm{m/s^2}$、
$R_E\simeq6.37\times10^6\,\mathrm m$、
$g_M\simeq1.62\,\mathrm{m/s^2}$、
$R_M\simeq1.737\times10^6\,\mathrm m$，
\[
\boxed{W_{\rm ext,min}\sim1.2\times10^{11}\,\mathrm J
\approx1.2\times10^8\,\mathrm{kJ}}.
\]
\excheck 这是一种准静态理想势能预算；
不能将其视为有时间限制、必须安全着陆的真实任务
的发动机燃料能量最小值。原书简答数量级约为
$1.3\times10^8\,\mathrm{kJ}$。

\exproblem{233}{36}
忽略地球自转，估算从地球表面发射物体，
使它最终能逃出\emph{整个太阳系}所需的最小初速，
即第三宇宙速度。
\exsolution
记地表第一宇宙速度
$v_{1,E}=\sqrt{GM_E/R_E}\simeq7.9\,\mathrm{km/s}$，
地球绕太阳的圆轨道速率
$V_E=\sqrt{GM_\odot/D}\simeq29.8\,\mathrm{km/s}$。
日地轨道半径处太阳逃逸速率为 $\sqrt2V_E$。
为使从地球引力井出来的物体在日心系中达到这一速度，
应沿地球公转方向发射，使地球外的双曲线余速
\[
u_\infty=(\sqrt2-1)V_E.
\]
地表初速 $v_3$ 的单位质量地球中心机械能为
$v_3^2/2-GM_E/R_E$，
在地球引力范围外为 $u_\infty^2/2$，
故
\[
\boxed{v_3=
\sqrt{\frac{2GM_E}{R_E}
+(\sqrt2-1)^2V_E^2}
=\sqrt{2v_{1,E}^2+(\sqrt2-1)^2V_E^2}
\simeq16.7\,\mathrm{km/s}}.
\]
\excheck 最优方向为地球公转切向前进方向；
忽略了地球自转、月球引力、其他行星摄动及真实轨道
交接处的细节，属于经典近似估算。
